\documentclass{beamer}

\usepackage[utf8]{inputenc}
\usepackage[T1]{fontenc}
\usepackage[french]{babel}
\usepackage{graphicx}
\usepackage{hyperref}

\usetheme{Madrid}
\usecolortheme{default}

\title[Projet SGSU]{Système de Gestion de la Scolarité Universitaire (SGSU)}
\subtitle{Présentation de Projet TP - Base de Données Avancée (BDDA)}
\author{NATOUM PIERRE}
\institute[BDDA]{Enseignant : BRAHIM ISSA HASSABALLAH}
\date{Année académique 2025-2026}

\begin{document}

\begin{frame}
    \titlepage
\end{frame}

\begin{frame}{Contexte et Objectifs}
    \textbf{Le problème actuel :} 
    \begin{itemize}
        \item Dans le système LMD, la gestion des dossiers étudiants, le calcul des moyennes et le suivi des paiements sont souvent chronophages.
        \item Sujet aux erreurs humaines sans système centralisé.
    \end{itemize}
    
    \vspace{0.5cm}
    \textbf{Notre solution (SGSU) :}
    \begin{itemize}
        \item Créer une plateforme web moderne et automatisée.
        \item Démontrer la modélisation d'une base de données relationnelle complexe.
        \item Connecter cette base à une application web fonctionnelle (TP BDDA).
    \end{itemize}
\end{frame}

\begin{frame}{Architecture Technique (La Stack)}
    Architecture cloud moderne séparant le Frontend et le Backend :
    
    \begin{itemize}
        \item \textbf{Frontend (Interface Utilisateur) :}
        \begin{itemize}
            \item React 18 et Vite.js
            \item Material UI (Design "Glassmorphism")
        \end{itemize}
        \item \textbf{Backend (Logique Métier) :}
        \begin{itemize}
            \item Serveur Node.js / Express (API REST)
        \end{itemize}
        \item \textbf{Base de données et Hébergement :}
        \begin{itemize}
            \item \textbf{PostgreSQL} hébergé dans le Cloud (Neon Serverless)
            \item Manipulée via l'ORM \textbf{Prisma}
            \item Déploiement automatisé sur \textbf{Render}
        \end{itemize}
    \end{itemize}
\end{frame}

\begin{frame}{Modèle Conceptuel des Données (MCD)}
    \textbf{Entités principales et associations :}
    \begin{itemize}
        \item \textbf{ETUDIANT} (1,N) $\longleftrightarrow$ \textit{S'inscrire} $\longleftrightarrow$ (1,N) \textbf{COURS}
        \item \textbf{INSCRIPTION} (1,N) $\longleftrightarrow$ \textit{Avoir\_Note} $\longleftrightarrow$ (1,1) \textbf{NOTE}
        \item \textbf{ETUDIANT} (0,N) $\longleftrightarrow$ \textit{Effectuer} $\longleftrightarrow$ (1,1) \textbf{PAIEMENT}
    \end{itemize}
    
    \vspace{0.3cm}
    \textbf{Règles de gestion :}
    \begin{itemize}
        \item Un étudiant peut s'inscrire à plusieurs cours, et un cours accueille plusieurs étudiants.
        \item Une inscription génère plusieurs notes (CC, Examen).
        \item Un paiement est obligatoirement rattaché à un seul étudiant.
    \end{itemize}
\end{frame}

\begin{frame}{Modèle Logique des Données (MLD)}
    Traduction relationnelle (Tables, Clés Primaires \underline{PK}, Clés Étrangères \textbf{\#FK}) :
    
    \begin{itemize}
        \item \textbf{Etudiant}(\underline{id}, matricule, nom, prenom, email, filiere, \dots)
        \item \textbf{Cours}(\underline{id}, code, titre, credits, semestre)
        \item \textbf{Inscription}(\underline{id}, \textbf{\#etudiantId}, \textbf{\#coursId}, date, statut)
        \item \textbf{Note}(\underline{id}, \textbf{\#inscriptionId}, type, note, coef, session)
        \item \textbf{Paiement}(\underline{id}, \textbf{\#etudiantId}, montant, trxId, methode, date)
    \end{itemize}
    
    \vspace{0.3cm}
    \textbf{Contraintes d'intégrité (SGBD) :} 
    \begin{itemize}
        \item Contrainte d'unicité (\texttt{UNIQUE}) sur \texttt{matricule} et le couple \texttt{(etudiantId, coursId)}.
        \item \texttt{ON DELETE CASCADE} pour garantir l'intégrité référentielle absolue.
    \end{itemize}
\end{frame}

\begin{frame}{Démonstration en Direct}
    \begin{center}
        \Large \textbf{Passons à la démonstration de l'application déployée !}
        
        \vspace{1cm}
        \normalsize Lien officiel de l'application :\\
        \url{https://sgsu-y5b2.onrender.com}
        
        \vspace{1cm}
        \normalsize Identifiants de test :\\
        Email : \texttt{admin@sgsu.edu}\\
        Mot de passe : \texttt{admin@123}
    \end{center}
\end{frame}

\begin{frame}{Bilan et Perspectives}
    \textbf{Conclusion :}
    \begin{itemize}
        \item Plateforme cloud opérationnelle.
        \item Preuve de la maîtrise de la chaîne complète : BD $\rightarrow$ Backend $\rightarrow$ Frontend $\rightarrow$ Hébergement.
    \end{itemize}
    
    \vspace{0.5cm}
    \textbf{Perspectives (V2.0) :}
    \begin{itemize}
        \item Migration du backend vers \textbf{Spring Boot (Java)}.
        \item Intégration de \textbf{Keycloak} pour sécuriser l'authentification.
        \item Génération automatique de Procès-Verbaux (PV) en PDF.
    \end{itemize}
\end{frame}

\begin{frame}
    \begin{center}
        \Huge \textbf{Merci de votre attention !}
    \end{center}
\end{frame}

\end{document}
